% Tabella 5.1 - confronto dei tre metodi proposti
\begin{table}[h]
\centering
\small
\renewcommand{\arraystretch}{1.35}
\setlength{\tabcolsep}{4pt}
\begin{tabular}{p{2.4cm} p{3.1cm} p{2.6cm} p{2.9cm} p{3.0cm}}
\toprule
\textbf{Metodo} & \textbf{Costo per iterazione} & \textbf{Memoria}
& \textbf{Ipotesi aggiuntive} & \textbf{Complessit\`a totale} \\
\midrule
Gradiente a campione dinamico &
$O(m\,n_k)$ &
$O(m)$ &
CCV soddisfatta a ogni iterazione; varianza limitata \eqref{eq:var_bound} &
$O\!\left(\dfrac{mL}{\lambda\varepsilon}\right)$ \\
Newton-CG con campionamento dinamico &
$O(m\,n_k + m\,n_h\,j_k)$\\($j_k$: iterazioni CG) &
$O(m\,n_h)$ (sottocampione Hessiana) &
Hessiana sdp sul sottospazio; criterio di arresto (5.35)--(5.36) &
$O\!\left(\dfrac{mL}{\lambda\varepsilon}\right)$ iter.\ esterne \\
Newton-CG con regolarizzazione $L_1$ &
come Newton-CG + costo active set &
$O(m)$ &
identificazione corretta della faccia ortante (5.41)--(5.46) &
come Newton-CG \\
\bottomrule
\end{tabular}
\caption{Confronto dei tre metodi proposti nel Capitolo~\ref{sec:algoritmi}.
$n_k$ \`e la dimensione del batch corrente, $n_h = R\,n_k$ quella del
sottocampione per l'Hessiana, $j_k$ il numero di iterazioni interne del
gradiente coniugato, $m$ il numero di variabili.}
\label{tab:5_1}
\end{table}
